\section{確率空間}
\subsection{可測空間}
%\begin{soudef}[$\sigma$加法族]
    $S$を空でない集合とする。$\mathcal{S}\subset 2^S$が$\sigma$加法族であるとは、
\begin{enumerate}
    \item $S\in \mathcal{S}$
    \item $A\in \mathcal{S}\Rightarrow A^c\in \mathcal{S}$
    \item $A_n\in \mathcal{S}\ \ (n\ge 1) \Rightarrow \bigcup_{n= 1}^\infty A_n\in \mathcal{S}$
\end{enumerate}
%\end{soudef}

%\begin{souexample}
    $2^S$は明らかに$\sigma$加法族
%\end{souexample}

%\begin{souprop}
    $S$を空でない集合、$\mathcal{S}$をその上の$\sigma$加法族とする。このとき
\begin{equation}
    A_n\in \mathcal{S}\ \ (n\ge 1)\Rightarrow \bigcap_{n= 1}^\infty A_n\in \mathcal{S}
\end{equation}
    が成立。
%\end{souprop}
%\begin{souremark}
    
%\end{souremark}
